\chapter{State Of Art}\label{sec:StateofArt}

The accurate estimation of battery internal states is an active research area driven by the growing 
adoption of electric vehicles, energy storage systems, and innovation in lithium-ion batteries. The literature
can be divided intro three approaches: the empirical, model-based and data-driven methods, each approach
differs in accuracy, computational cost, and implementation complexity~\cite{SOHBook}. While reviewing estimation approaches,
it is necessary to also consider the battery models used to describe the battery's dynamics.

\subsection{Empirical Estimation Methods}

One of the simplest models for \acs{SOC} estimation is the Coulomb Counting (\acs{CC}), also known as Ampere-hour integration method.
This is an open-loop approach that computes the State of Charge by integrating the measured cell current over time. 
The advantages of this method rely on its simple implementation and low computational demand~\cite{BatteryKalmanRaceCar}, 
nonetheless, its performance is highly dependent on the initial \acs{SOC} characterisation, furthermore \acs{CC} is susceptible 
to sensor inaccuracies and accumulated integration errors~\cite{EKFEV}\cite{SOCSOH}. 
\\
\\
A complementary approach to \acs{CC} is the Open-Circuit Voltage (\acs{OCV}) method, which applies the relationship between \acs{OCV} and \acs{SOC}
to infer the battery state from a voltage measurement. This \acs{OCV}-\acs{SOC} curve is commonly characterised experimentally and 
stored as a lookup table~\cite{SOHBook}. Despite its simplicity and applicability across different battery chemistries, 
the \acs{OCV}-\acs{SOC} relationship only holds accurately when the cell is fully relaxed and undisturbed for several hours p
rior to measurement. This condition can only be achieved under controlled laboratory conditions, making the method unsuitable 
for applications involving dynamic loads~\cite{plett2015battery}.
As a result, both the \acs{CC} and \acs{OCV} are typically integrated into more complex estimation methods than used independently.

\subsection{Battery Modelling Approaches}\label{SoABatteryModels}

To review model-based estimation techniques, different types of battery models need to be considered first, 
as the choice of model conditions the achievable accuracy and the feasibility of the implementation.
\\
\\
In the literature there are two different methods to model a battery, the Physics-based models (\acs{PBM}) and the Equivalent 
Circuit models (\acs{ECM})~\cite{plett2015battery}.~\acsp{PBM} capture the internal dynamics of a battery cell by constructing the 
equations from electrochemical and thermodynamic laws, which describe the physical transport and diffusion and 
electrochemical processes occurring in the cell. This type of model offers high predictive accuracy across a wide 
range of operating conditions and insights into the internal electrochemical states, useful for predicting ageing. 
However, the mathematical complexity involves coupled partial differential equations and demands high computational
effort, which results in high costs and multiple challenges to implement on an embedded \acs{BMS}~\cite{plett2015battery}.
\\
\\
\acsp{ECM} represent the battery cell as an electric circuit composed of resistors, capacitors, and an ideal voltage source, providing
an approximate behaviour of the current-voltage dynamics. Its accuracy is generally lower than that of physics-based models, 
and its predictive validity is restricted to the specific operational conditions in which it was characterised~\cite{plett2015battery}
~\cite{SOHBook}.
However, its simple structure, parameterisation process, and suitability to implement in real applications offer significant 
practical advantages. Although the parameterisation process demands systematic and methodical laboratory testing, 
the model's low mathematical complexity makes it more feasible to implement in comparison to \acsp{PBM}.
The parameter values of an \acs{ECM} depend on operating conditions such as temperature, C-rate, depth of discharge, 
and ageing status, which must be accounted for in the model structure~\cite{DEVARAKONDA2014}. Among \acs{ECM} variants, the second-order 
Thévenin model has emerged as the standard choice in the literature. This circuit comprises of an \acs{OCV} source in series with a 
resistor connected in series with two parallel branches, each with a capacitor and a resistor. 
This model captures both electrochemical polarisation and concentration polarisation dynamics while 
remaining computationally advantageous for \acs{BMS} applications~\cite{BatteryKalmanRaceCar}\cite{EKFEV}\cite{SOCSOH}.

\subsection{Adaptive and Fuzzy Logic Methods for Estimation}

To improve on fixed-parameter limitations of the standard model-based estimators, works have proposed adaptive strategies
that update model parameters. The main challenge is the online parameter identification, traditional offline methods based on least-squares 
curve fitting introduce accuracy limitations, while standard online approaches such as the Recursive Least Squares (\acs{RLS}) 
algorithm suffer from data saturation and sensitivity to initial value settings, making it difficult 
to reliably track time-varying parameters~\cite{Wang2025LithiumIonBS}.
Zhang et al.~\cite{UKFDynamPar} demonstrated that by incorporating dynamic, \acs{SOC}-dependent parameters identified through limited-memory 
recursive least-squares algorithm, the results show a more representative model of the actual battery discharge behaviour, with
the resulting \acs{UKF}-based \acs{SOC} estimates closely matching those obtained via the ampere-hour method used as a reference benchmark.
\\
\\
A further extension of model-based estimation involves the integration of fuzzy logic to dynamically adjust 
filter parameters during operation. Wang et al.~\cite{Wang2025LithiumIonBS} proposed combining the 
Multiple Innovation Least Squares (\acs{MILS}) algorithm for online parameter 
identification with a fuzzy-logic mechanism that continuously adapts the observation noise covariance matrix of an \acs{EKF}. 
The \acs{MILS} algorithm proves to enhance the data efficiency and convergence speed while the fuzzy-logic optimises the balance
between the time update and measurement correction of the filter to improve higher precision under the coupled and dynamic conditions 
of an electric vehicle. This type of hybrid methods represents a more complex implementation, but addresses some of 
the robustness limitations of the Kalman filter approaches.

\subsection{Data-Driven and Machine-Learning Models}

Growing interest in data-driven methods has risen as an alternative to model-based approaches, due to their effectiveness in handling
nonlinearities without requiring characterising the internal battery dynamics. These methods require experimental data
when training to learn the input/output relationship~\cite{SOHBook}. Neural-Network estimators, such as the backpropagation (\acs{BP})
have been applied to \acs{SOC} across different battery chemistries and operating conditions~\cite{SOCLowSpeedNeuralNetwork}.
Xu et al. demonstrated that for low-speed electric vehicle batteries, \acs{BP} neural networks that have
sufficiently deep network topologies can achieve suitable \acs{SOC} prediction accuracy during charging cycles, with the understanding that
results are sensitive to the network architecture and the size and diversity of the training dataset~\cite{SOCLowSpeedNeuralNetwork}.
\\
\\
In Bosello's work~\cite{Bosello2025} different deep learning approaches are provided for \acs{Li-ion} battery state estimation, 
demonstrating the competitive accuracy across a range of conditions. Nonetheless, the consistent limitation in all data-driven 
methods is their dependency on large and high quality training datasets that must represent the operating conditions faced~\cite{EKFEV}
~\cite{Bosello2025}. Consequently, for applications where training data for the specific cell under its operating conditions is
insufficient, model-based approaches are more practical.

\subsection{Kalman Filter Methods for \acs{SOC} and \acs{SOH} Estimation}

Kalman filter algorithms have demonstrated consistency, feasibility, and cost-effective performance as battery estimators~\cite{EKFEV}.
The Linear Kalman Filter is optimal for linear systems with Gaussian noise but cannot directly accommodate the nonlinearities of 
a lithium-ion \acs{OCV}/\acs{SOC} relationship~\cite{SOCSOH}. The Extended Kalman Filter addresses this by linearising the nonlinear model at each
time step through a first-order Taylor expansion, followed by the standard recursive Kalman filter~\cite{plett2015battery}.
\\
\\
Luo et al.~implemented an \acs{EKF} combined with a second-order RC \acs{ECM} for \acs{Li-ion} batteries on an electric vehicle, validating the algorithm
under different current profiles: under constant current discharge with a maximum estimation error of $0.9\%$, under constant pulse and 
periodic variable current discharges where the error stayed within $1.2\%$ and $1.8\%$ respectively, and under rapidly varying current
discharge, obtaining an error of $3.4\%$~\cite{EKFEV}. Gjengedal et al.~applied a comparable \acs{EKF} approach to a Formula Student 
electric race car, using lithium-polymer cells which were subjected to high and rapidly varying current profiles. 
Their implementation achieved a \acs{SOC} estimation error below $0.66\%$, proving that Kalman filter estimators
are a suitable option for motorsport applications~\cite{BatteryKalmanRaceCar}.
\\
\\
The Sigma-Point or Unscented Kalman Filter (\acs{UKF}) avoids the approximation error introduced by the \acs{EKF}'s linearisation 
step by performing a statistical linearisation at each time step rather than computing Jacobians~\cite{plett2015battery}. 
It is concluded that the \acs{UKF} avoids the linearisation error inherent to the \acs{EKF} and achieves higher estimation precision, at 
the cost of increased computational complexity.
\\
\\
For simultaneous estimation of the \acs{SOC} and \acs{SOH} parameters, the joint and dual estimation techniques have been developed.
Liu et al.~have proposed a joint online \acs{SOC} and \acs{SOH} estimation method using a Dual \acs{EKF}, in which one filter estimates \acs{SOC} at a fast 
sampling period of 1 second while a second filter estimates the \acs{SOH}-related parameters at 60-second intervals~\cite{SOCSOH}. 
The method was tested using a second-order RC model of a ternary lithium cell under both constant current and random current 
profiles, achieving \acs{SOC} and \acs{SOH} estimation accuracy that outperformed standalone \acs{EKF} and \acs{UKF} implementations.

\subsection{Contribution}

The reviewed literature demonstrates that an Extended Kalman Filter approach combined with a second-order RC 
equivalent circuit model provides a validated and feasible method to estimate the State of Charge of a battery 
under dynamic conditions. Incorporating \acs{UKF}, adaptive parameter identification, and dual \acs{SOC}/\acs{SOH} estimation have 
further improved the accuracy and robustness of the estimation. 
However, these adjustments have largely been studied in isolation from one another and from the specific demands of motorsport applications. 
Gjengedal et al.~\cite{BatteryKalmanRaceCar} validated an \acs{EKF} \acs{SOC} estimator on a Lithium Polymer (\acs{Li-Po}) cell under race 
current profiles, but 
did not extend the framework to \acs{SOH} estimation. Similarly, Liu et al.~\cite{SOCSOH} demonstrated a Dual \acs{EKF} approach for 
joint \acs{SOC} and \acs{SOH} estimation on a ternary lithium cell for stationary energy storage and not for a motorsport application.
Furthermore, conventional parameter identification methods used across the reviewed studies require extended rest periods 
between tests, a limitation that the algebraic method of Devarakonda et al.~\cite{DEVARAKONDA2014} is specifically designed to overcome.
\\
This work combines the experimental characterisation of the WeLion SHP270\-30 cell, efficient parameter identification 
through the algebraic method, and a Dual \acs{EKF} approach for joint \acs{SOC} and \acs{SOH} estimation, intended for an eRacing 
electric motorcycle running under dynamic operating conditions.